\renewcommand{\abstractnamefont}{\normalfont}
\renewcommand{\abstractname}{\uppercase{Abstract}}
\begin{abstract}
  \begin{center}
    \thetitle\\
    \theauthor
  \end{center}
  \vspace{1cm}
\pagestyle{empty}

Algebraic Geometry was revolutionized in the 1960s by many incredible techniques introduced by Grothendieck and his collaborators. Arguably 
the most fundamental machinery was the introduction of sheaf cohomology, which in turn relied on the developments of homotopical algebra. 
Here, the notion of flatness becomes a crucial technical condition, especially when considering the interaction between base-change and cohomology, where the flatness condition 
is imposed on a morphism of schemes $f: X \to Y$ to rid $X$ of any pathological behavior for the resulting family of fibers 
$X_y$ for $y \in Y$. It would be unreasonable, however, to expect that the flatness condition on $f$ is generally sufficient to guarantee good behavior of the family of fibers $X_y$, 
because flatness is a module-theoretic condition that can generally exhibit quite strange behavior without additional finiteness hypotheses. Broadly speaking,
the pursuit of understanding this dichotomy is an underlying theme for this thesis.\\
\indent Semirings, studied in Chapter 1, provided new aspects of this dichotomy to wrestle with. One of our main results is that 
the most natural notion of homotopy theory for semirings is not enough to capture the properties of flatness required to guarantee good base-change properties.
\begin{theorem}[See Corollary~\ref{cor:main}]
For any commutative monoid $M \in \mod{\bb{N}}$, $M \otimesl_{\bb{N}} \bb{Z}$ is a discrete module. On the other hand, 
the semiring map $\bb{N} \to \bb{Z}$ is not flat, in the sense that the functor $-\otimes_{\bb{N}} \bb{Z}:\mod{\bb{N}} \to \mod{\bb{N}}$ does not preserve finite limits (see Corollary~\ref{cor:NRnotflat}).
\end{theorem}
While this result is indeed in stark contrast to what happens in classical commutative algebra, where flatness of a ring map $R \to S$ can be detected via the derived functors of $-\otimesl_R S$, 
the semiring situation does bear some similarities to the classical situation.
\begin{theorem}[See Theorem~\ref{theorem:main}]
A flat, finitely presented epimorphism of semirings $f: A \to B$ induces an isomorphism of $\text{Spec}(B)$ onto a quasi-compact open of $\text{Spec}(A)$ in the Zariski-topology.
\end{theorem}
We were informed after having proved Theorem~\ref{theorem:main} that Florian Murty obtained the same result in greater generality via
a different method (see \cite{Martyflat}).\\
\indent Our main interest in studying the flatness of semirings is that here the notion seems to capture some strong positivity properties of semirings. The following is 
a question posed to the author by James Borger.
\begin{question}
Let $A$ be a finitely presented flat $\bb{R}_{+}$-algebra. 
If $A$ is non-trivial, does it necessarily have an $\bb{R}$-point (or even better, an $\bb{R}_{+}$-point)?
\end{question}
If we asked the above question for the groupification $\bb{R}$ of $\bb{R}_{+}$, then the answer would be no, as there are many non-trivial flat and finitely presented $\bb{R}$-algebras that do not have any $\bb{R}$-points.
In our opinion, the question seems to suggest that because $\bb{R}_{+}$ is a totally ordered semiring with no elements admitting additive inverses, the flatness condition on $A$ is strong enough to guarantee the existence of a $\bb{R}_{+}$-point.
We attempted to address this question by understanding when quotients of the form $\bb{R}_{+}[x_1,...,x_n]/(f \sim g)$ are flat over $\bb{R}_{+}$, but we were not able to obtain 
much progress. In private communication, Johan de Jong informed the author of the following result.
\begin{theorem}[de Jong]
The semiring $\bb{R}_{+}[x]/(ax \sim x^2 + 1)$ is flat over $\bb{R}_{+}$ if and only if $a \ge 2$. 
\end{theorem}

On the other hand, in Chapter 2, we will show that from the perspective of \textit{descendability}, flat ring maps between classical rings can behave 
arbitrarily badly. The notion of a descendable morphism is a technical condition introduced by Akhil Mathew \cite{DescendMathew} 
building upon the work of Paul Balmer \cite{DescendBalmer}. Informally, if a map $f: A \to B$ is descendable, then $A$-linear stable $\infty$-categories can
be understood via $B$-linear stable $\infty$-categories together with `descent' data (see \ref{thm:luriedescend}). Our main result is that 
faithfully flat ring maps are not necessarily descendable; see Corollary~\ref{corollary:mainindiv}, Theorem~\ref{theorem:maincounter}, and the work of Aoki \cite{aoki}. 
This answers a question by Bhatt and Scholze (\cite[11.24]{BhattScholzeProj}), Akhil Mathew (\cite[Prop. 3.32]{DescendMathew}), and Jacob Lurie (\cite[D.3.3.4]{LurieSAG}). 
Our strategy was to first understand the relationship between non-vanishing cup-products and cardinality in the module-theoretic setting.
\begin{theorem}[See Theorem~\ref{theorem:indivisiblecup}]
For a commutative ring $R$, if
$R$ contains an $n$-indivisible sequence (see Definition~\ref{definition:indivisible}), then 
there exists a flat module $M$ over $R$ and a class $\eta \in \Ext^1_R(M, R')$ such that 
\[\eta^{\otimesl_R n} \in \Ext^n_R(M^{\otimesl_R n}, R'^{\otimesl_R n}) \neq 0\]
where $R'$ is a free $R$-module.
\end{theorem}
We then form a quite general method to convert module theoretic non-vanishing cup-products to non-descendable ring maps, and consequently obtain many examples of non-descendable faithfully flat ring maps (See Corollary~\ref{corollary:mainindiv}).
Our method in particular allows us to construct examples between $p$-boolean rings, which was a question posed to the author by Juan Esteban Rodríguez Camargo. By slightly refining our argument, we were also able to construct a faithfully flat cover of $\spec{k[x_1,x_2,...]}$, where $k$ is an algebraically closed field, that is not descendable (see Theorem~\ref{theorem:maincounter}).
This example is quite striking, as $\spec{k[x_1,x_2,...]}$, while not Noetherian, is still a reasonably nice scheme from the perspective of algebraic geometry.\\

To wrap up this thesis, we homed in on a conjecture posed by Grothendieck and Dieudonné. In \cite[21.12.14]{EGAIV}, 
the authors conjectured that for a locally of finite type map $f: X \to Y$ of excellent, locally Noetherian schemes, with $X$ normal and $Y$ regular, the 
ramification locus of $f$ is pure of codimension $1$ (See Theorem~\ref{thm:nagata-purity}). Having already shown that for any open subscheme $V \subset X$ such that $V \to X$ is an affine morphism of schemes,
 $X \setminus V$ is pure of codimension $1$ (\cite[21.12.7]{EGAIV}), the authors then conjectured that
if $V$ is the maximal open subscheme $V \subset X$ where $f$ is unramified (and hence \'{e}tale), then $V \to X$ is an affine morphism of schemes (see (v) \cite[21.12.14]{EGAIV}).
The main goal of this chapter was to answer their conjecture in the affirmative, see Theorem~\ref{thm:affine-mixed}. We will in particular demonstrate the following:
\begin{theorem}[See Theorem~\ref{thm:cohpuregen}]
Let $(A,\ideal{m}_A,k)$ be a regular local ring and $V \to \spec{A}$ an \'{e}tale morphism that is cohomologically pure in codimension $1$. Then $V$ is an affine scheme.
\end{theorem}
Our definition of `cohomologically pure in codimension $1$' (see Definition~\ref{def:cohpure}) could be seen as a cohomological analogue of the situation that $V$ arises 
as an open subscheme of $\spec{B}$, for a normal local ring $(B, \ideal{m}_B)$ that is finite over $A$, such that there is a factoring
$ V\to \spec{B} \setminus \{\ideal{m}_B\} \to \spec{B}$ where the first map is affine morphism of schemes. We only realized later that the authors in \cite[21.12]{EGAIV} were interested in characterizing rings $B$ for which such a $V$ is then 
forced to be an affine scheme, and consequently formulated a conjecture (iv) \cite[21.12.14]{EGAIV}.\\
\indent If their conjecture was true, 
then as noted in (v) \cite[21.12.14]{EGAIV}, the affineness of the maximal \'{e}tale locus would be equivalent to the purity of the ramification locus. 
By essentially paraphrasing the proof of purity of the ramification locus given in \cite[0ECD]{sp}, 
Theorem~\ref{thm:affine-mixed} would then follow by induction on the dimension of $Y$, where the case of dimension $2$ would crucially use 
the \textit{miracle flatness theorem} (see \cite[00R4]{sp}), from which this thesis derives its name.
\end{abstract}
\thispagestyle{empty}